\documentclass[../otf-unofficial-guide]{subfiles}
\begin{document}

\section{カウンタに\texorpdfstring{\cmd{aj}}{aj}系囲みつき文字を利用する}

本ドキュメントでも実践しているように、ページ番号や脚注・傍注の番号、enumerate環境の番号等を\pkg{ajmacros}から提供される\cmd{aj}系囲みつき文字コマンドに変更できます。

例えば、\cmd{ajMaru}を使って丸囲み数字にする場合、\cmd{ajMaru\{10\}}とすれば``\nobreak\ajMaru{10}\nobreak''となります。あるいは、\cmd{arabic}のようにカウンタ名を引数としたい場合は、\cmd{ajLabel}\cmd{ajMaru\{\textlrangle{counter}\}}とします。
\sidenote{\cmd{ajLabel}は引数をカウンタに変更するためのコマンドですが、\pkg{hyperref}の更新によってはエラーが生じることがあります。\\{\scriptsize \textgt{参考}：\href{https://note.golden-lucky.net/2016/02/latexajmaruhyperref.html}{k16's note: \LaTeX で連番リストに丸数字（\cmd{ajMaru}）を使えたはずだけどhyperrefで困る話}}}

\bigskip

より具体的に、本ドキュメントでの実例を紹介します。

\texttt{page}のラベルを\pkg{ajmacros}から提供される\cmd{ajKakkoKansuji}コマンドを利用して括弧囲み漢数字に変更する場合は以下のようにします。

\begin{texsource}
\renewcommand{\thepage}{\ajLabel\ajKakkoKansuji{page}}
\end{texsource}

あるいは、\pkg+{enumitem}を使用する場合は、以下のように\cmd{protect}を使用します。

\begin{texsource}
\begin{enumerate}[label=\protect\ajMaru{\arabic{*}}]
  ...
\end{enumerate}
\end{texsource}

\cmd{ajMaru}等のコマンドはfragileなコマンドのため、\cmd{section}内であっても\cmd{protect}を利用する必要があります。

\begin{texsource}
\section{\protect\ajMaru{5}}
\end{texsource}




\end{document}
